% Standalone gallery for manuscript pentagon_v005.tex.
%%
% Author Thomas C. Hales et al.
% LaTeX Format


%!TEX TS-program = latex    
%% This line is for TexShop. 

% started June 2015
% submitted to arxiv Feb 23, 2016

\documentclass{amsart}
\usepackage{graphicx}
\usepackage{amsfonts}
\usepackage{amscd}
\usepackage{amssymb}
\usepackage{alltt}
\usepackage{multicol}
\usepackage{amsmath}
\usepackage{amsthm}
\usepackage{amscd}
%\usepackage{mnsymbol}

%%\documentclass[spanningrule]{cambridge7A}
%% Cambridge University Press Macros from
%% https://authornet.cambridge.org/information/productionguide/laTex_files/


% required by CUP
\usepackage[numbers]{natbib}
\usepackage{rotating}
\usepackage{floatpag}
 \rotfloatpagestyle{empty}
\usepackage{amsthm}
\usepackage{graphicx}
%\usepackage{multind}\ProvidesPackage{multind}
\usepackage{times}


% my additions
\usepackage{verbatim}
\usepackage{latexsym}
\usepackage{amsfonts}
\usepackage{amsmath}
\usepackage{crop}
\usepackage{fancyhdr}
\pagestyle{fancy}
\usepackage{txfonts}
%\usepackage[hyphens]{url}
%\usepackage{url}
\usepackage{setspace}
\usepackage{ellipsis} % 
% http://www.ctan.org/tex-archive/macros/latex/contrib/ellipsis/ellipsis.pdf 
% wrapping graphics
% http://en.wikibooks.org/wiki/LaTeX/Floats,_Figures_and_Captions
\usepackage{wrapfig}
\usepackage{listings}


%tikz graphics
\usepackage{xcolor} % to remove color.
\usepackage{tikz} % Needs pgf version >= 2.10.
\usetikzlibrary{chains,shapes,arrows,trees,matrix,positioning,backgrounds,fit,calc,fadings}

% fonts
\usepackage[mathscr]{euscript} % powerset.
\usepackage{pifont} %ding
\usepackage[displaymath]{lineno}


%%%%%%%%%%%%%%%%%%%%%%%%%%%%%%%%%%%%%%%%%
\def\tikzfig#1#2#3{%
\begin{figure}[htb]%
  \centering
\begin{tikzpicture}#3
\end{tikzpicture}
  \caption{#2}
  \label{fig:#1}%
\end{figure}%
}

\def\tikzwrap#1#2#3#4{%
\begin{wrapfigure}{r}{#4\textwidth}
  \begin{center}
\begin{tikzpicture}#3
\end{tikzpicture}
\end{center}
  \caption{#2}
  \label{fig:#1}%
\end{wrapfigure}%
}

%-%
%--Redefining--
%-%
\def\emptyset{\varnothing}
\def\ups{\upsilonup} % Needs txfonts; else use \upsilon

%% % vector notation
\def\v{{\mathbf v}}
\def\u{{\mathbf u}}
\def\w{{\mathbf w}}
\def\p{{\mathbf p}}
\def\q{{\mathbf q}}

% operatorname
\def\op#1{{\operatorname{#1}}}
\def\optt#1{{\operatorname{{\texttt{#1}}}}}
\def\card{\op{card}}
\def\arc{\operatorname{arc}}
\newcommand{\ring}[1]{\mathbb{#1}}
\def\norm#1#2{\hbox{\ensuremath{\|#1\unskip-\unskip{#2}\|}}}
\def\normo#1{{\|#1\|}}

\def\red#1{\textcolor{red}{#1}}
\def\calc{$^\dagger$}
\def\calcref#1{$^\dagger$}
% line numbers
\def\lll{\resetlinenumber[1]}
\def\linenumberfont{\normalfont\small\sffamily}

\def\tocpart#1{
  \addcontentsline{toc}{part}{\Large{#1}}}
% (or even \LARGE)
  
\crop
%\makeindex
%\makeindex{index/Notation}
%\makeindex{index/Index}

\def\linput#1{\lll\input{#1}}


\raggedbottom  % for now.
%\raggedright  % don't worry for now.

%% \renewcommand{\baselinestretch}{2}  %% turns on double spacking

%%% end notes -kill this.
%%\maketextnotes

%%%%%%%%%%%%%%%%%%%%%%%%%%%%%%%%%%%%%%%%

% new theorems

% One section-based counter shared by equations and theorem-like environments.
\numberwithin{equation}{section}
% Figures and tables have their own section-based sequences.
\numberwithin{figure}{section}
\numberwithin{table}{section}

\theoremstyle{plain}
\newtheorem{theorem}[equation]{Theorem}
%\newtheorem{theorem*}[equation]{Theorem$^*$}
\newtheorem{lemma}[equation]{Lemma}
%\newtheorem{lemdag}[equation]{Lemma$^\dagger$}
\newtheorem{background}[equation]{Background}
\newtheorem{corollary}[equation]{Corollary}
\newtheorem{example}[equation]{Example}
\newtheorem{conjecture}[equation]{Conjecture}
%---
\theoremstyle{definition}
\newtheorem{definition}[equation]{Definition}
%---
\theoremstyle{remark}
\newtheorem{remark}[equation]{Remark}
\newtheorem{notation}[equation]{Notation}
\newtheorem{exer}[equation]{Exercise}

\renewcommand\footnotemark{}
%\renewcommand\footnoterule{}

\def\fig#1#2#3{%
\begin{figure}[htb]%
  \centering
#3
  \caption{#2}
  \label{fig:#1}%
\end{figure}%
}


%%%%%%%%%%%%%%%%%%%%%%%%%%%%%%%%%%%%%%%%%


\usepackage[margin=24mm]{geometry}
\begin{document}
% Extracted manuscript macros; gallery-only label display overrides below.
\def\area{\op{area}}
\def\areta{\op{area}^\circ}
\def\ao{a_\text{min}}
\def\acrit{a_\text{crit}}
\def\loc{\op{loc}}

\def\Ra{\Rightarrow}
\def\nRa{\nRightarrow}
\def\rab{\Ra_{b}}

\def\C{\mathcal C}
\def\SS{\mathcal S}
\def\N{\mathcal N}
\def\M{\mathcal M}
\def\T{\mathcal T}
\def\PD{\Psi D}
\def\cong{\equiv}
\def\r{{\mathbf r}}
\def\s{{\mathbf s}}
\def\bl{\ell}
\def\epsM{\epsilon_{\M}}
\def\epsN{\epsilon_{\N}}
\def\c{{\mathbf c}}
\def\K{\op{K}}


\def\dx#1{d_{#1}} 
\def\dd#1#2{\norm{\c_{#1}}{\c_{#2}}}
\def\tb{\underline{b}}
\def\tn{\bar n}

%for debugging:
\def\libel#1{{\text{\sc [#1]~}}\label{#1}}
\def\rif#1{(\ref{#1}-{\text{\sc #1})}}
%\def\libel#1{\label{#1}}
%\def\rif#1{\ref{#1}}

% tikz.
\def\smalldot#1{\draw[fill=black] (#1) node [inner sep=0.8pt,shape=circle,fill=black] {}}
\def\graydot#1{\draw[fill=gray] (#1) node [inner sep=1.3pt,shape=circle,fill=gray] {}}
\def\whitedot#1{\draw[fill=gray] (#1) node [inner sep=1.3pt,shape=circle,fill=white,draw=black] {}}
\tikzset{dartstyle/.style={fill=black,rotate=-90,inner sep=0.7pt,dart,shape border uses incircle}}
\tikzset{grayfatpath/.style={line width=1ex,line cap=round,line join=round,draw=gray}}

% parameters (x,y) center coords, theta vertex angle, rho radius in cm.
\def\pent#1#2#3#4{%
\draw[red] (#1,#2) + (#3:#4cm) -- + (#3+72:#4cm) -- 
+(#3+144:#4cm) -- +(#3+216:#4cm) -- + (#3+288:#4cm) -- cycle
}

\def\pentpink#1#2#3#4#5{%
\draw[#5] (#1,#2) + (#3:#4cm) -- + (#3+72:#4cm) -- 
+(#3+144:#4cm) -- +(#3+216:#4cm) -- + (#3+288:#4cm) -- cycle
}

% parameters (x,y) center coords, theta vertex angle.
\def\pen#1#2#3{%
\draw[red] (#1,#2) + (#3:1cm) -- + (#3+72:1cm) -- 
  +(#3+144:1cm) -- +(#3+216:1cm) -- + (#3+288:1cm) -- cycle
}

% parameters (x,y) center coords, theta vertex angle.
\def\penp#1#2{%
\draw[red] (#1) + (#2:1cm) -- + (#2+72:1cm) -- 
  +(#2+144:1cm) -- +(#2+216:1cm) -- + (#2+288:1cm) -- cycle
}


\def\threepent#1#2#3#4#5#6#7#8#9{%
\pen{#1}{#2}{#3};
\pen{#4}{#5}{#6};
\pen{#7}{#8}{#9};
\draw[blue] (#1,#2) -- (#4,#5) -- (#7,#8) -- cycle
}

\def\threepentnoD#1#2#3#4#5#6#7#8#9{%
\pen{#1}{#2}{#3};
\pen{#4}{#5}{#6};
\pen{#7}{#8}{#9}
}

\def\blankfig#1{\tikzfig{#1}{Insert graphic}
{
[scale=1]
\draw(0,) circle(1cm);
}}





\def\libel#1{\label{#1}}
\def\rif#1{\ref{#1}}

\makeatletter
\expandafter\def\csname r@sec:intro\endcsname{{1}{2}}
\expandafter\def\csname r@thm:main\endcsname{{1.1}{2}}
\expandafter\def\csname r@fig:double-lattice\endcsname{{1.1}{3}}
\expandafter\def\csname r@sec:cluster\endcsname{{2}{6}}
\expandafter\def\csname r@eqn:epsN\endcsname{{2.1}{6}}
\expandafter\def\csname r@eqn:epsM\endcsname{{2.2}{6}}
\expandafter\def\csname r@fig:constant\endcsname{{2.1}{6}}
\expandafter\def\csname r@fig:KKtriangle\endcsname{{2.2}{7}}
\expandafter\def\csname r@fig:subcritical\endcsname{{2.3}{7}}
\expandafter\def\csname r@fig:KK2min\endcsname{{2.4}{7}}
\expandafter\def\csname r@fig:Gamma\endcsname{{2.5}{8}}
\expandafter\def\csname r@fig:pseudo-dimer\endcsname{{2.6}{8}}
\expandafter\def\csname r@fig:obtuse\endcsname{{2.7}{9}}
\expandafter\def\csname r@fig:visual\endcsname{{2.8}{11}}
\expandafter\def\csname r@lemma:M-nonobtuse\endcsname{{2.7}{12}}
\expandafter\def\csname r@lemma:obtuse-b\endcsname{{2.8}{12}}
\expandafter\def\csname r@cor:obtuse-b-arrow\endcsname{{2.9}{12}}
\expandafter\def\csname r@lemma:n1m1\endcsname{{2.10}{12}}
\expandafter\def\csname r@lemma:n2m2\endcsname{{2.11}{12}}
\expandafter\def\csname r@lemma:main\endcsname{{2.12}{12}}
\expandafter\def\csname r@eqn:main\endcsname{{2.13}{12}}
\expandafter\def\csname r@eqn:strict-main\endcsname{{2.14}{13}}
\expandafter\def\csname r@thm2:main\endcsname{{2.15}{13}}
\expandafter\def\csname r@cor:equal\endcsname{{2.16}{14}}
\expandafter\def\csname r@sec:contact\endcsname{{3}{14}}
\expandafter\def\csname r@fig:receptor\endcsname{{3.1}{15}}
\expandafter\def\csname r@fig:3C-type\endcsname{{3.2}{16}}
\expandafter\def\csname r@lemma:sep1\endcsname{{3.1}{16}}
\expandafter\def\csname r@fig:clover\endcsname{{3.3}{17}}
\expandafter\def\csname r@lemma:sep2\endcsname{{3.2}{17}}
\expandafter\def\csname r@lemma:sep3\endcsname{{3.3}{17}}
\expandafter\def\csname r@fig:sep\endcsname{{3.4}{17}}
\expandafter\def\csname r@fig:delta\endcsname{{3.5}{18}}
\expandafter\def\csname r@fig:pin1\endcsname{{3.6}{18}}
\expandafter\def\csname r@sec:area\endcsname{{4}{19}}
\expandafter\def\csname r@lemma:21\endcsname{{4.1}{20}}
\expandafter\def\csname r@lemma:right\endcsname{{4.2}{20}}
\expandafter\def\csname r@lemma:primary\endcsname{{4.5}{21}}
\expandafter\def\csname r@lemma:primary2\endcsname{{4.6}{21}}
\expandafter\def\csname r@lemma:delta\endcsname{{4.7}{21}}
\expandafter\def\csname r@lemma:mid-172\endcsname{{4.8}{22}}
\expandafter\def\csname r@lemma:172\endcsname{{4.9}{22}}
\expandafter\def\csname r@fig:172-slider\endcsname{{4.1}{23}}
\expandafter\def\csname r@lemma:2C\endcsname{{4.10}{23}}
\expandafter\def\csname r@lemma:a0\endcsname{{4.11}{24}}
\expandafter\def\csname r@sec:calc\endcsname{{5}{25}}
\expandafter\def\csname r@NKQNXUN\endcsname{{5.1}{25}}
\expandafter\def\csname r@calc:pseudo1\endcsname{{5.1}{25}}
\expandafter\def\csname r@RWWHLQT\endcsname{{5.2}{25}}
\expandafter\def\csname r@calc:pseudo2\endcsname{{5.2}{25}}
\expandafter\def\csname r@calc:pseudo3\endcsname{{5.3}{25}}
\expandafter\def\csname r@BXZBPJW\endcsname{{5.4}{25}}
\expandafter\def\csname r@calc:pseudo-area\endcsname{{5.4}{25}}
\expandafter\def\csname r@JQMRXTH\endcsname{{5.5}{25}}
\expandafter\def\csname r@calc:pseudo-area3\endcsname{{5.5}{25}}
\expandafter\def\csname r@lemma:egress'\endcsname{{5.6}{25}}
\expandafter\def\csname r@lemma:m2-area\endcsname{{5.7}{25}}
\expandafter\def\csname r@KUGAKIK\endcsname{{5.10}{26}}
\expandafter\def\csname r@calc:dimer-isosc\endcsname{{5.10}{26}}
\expandafter\def\csname r@HUQEJAT\endcsname{{5.11}{26}}
\expandafter\def\csname r@calc:pent3\endcsname{{5.11}{26}}
\expandafter\def\csname r@QPJDYDB\endcsname{{5.12}{26}}
\expandafter\def\csname r@calc:pent4\endcsname{{5.12}{26}}
\expandafter\def\csname r@def:large\endcsname{{5.13}{27}}
\expandafter\def\csname r@fig:large\endcsname{{5.1}{27}}
\expandafter\def\csname r@FHBGHHY\endcsname{{5.14}{27}}
\expandafter\def\csname r@calc:large\endcsname{{5.14}{27}}
\expandafter\def\csname r@lemma:large\endcsname{{5.15}{27}}
\expandafter\def\csname r@sec:dimer-pair\endcsname{{6}{27}}
\expandafter\def\csname r@thm:dimer-area\endcsname{{6.1}{27}}
\expandafter\def\csname r@sub:long\endcsname{{6.1}{28}}
\expandafter\def\csname r@sec:O2C\endcsname{{6.2}{29}}
\expandafter\def\csname r@fig:slider\endcsname{{6.1}{30}}
\expandafter\def\csname r@sec:squeeze\endcsname{{6.3}{32}}
\expandafter\def\csname r@lemma:squeeze-area\endcsname{{6.5}{33}}
\expandafter\def\csname r@lemma:dimer-acrit-epsN\endcsname{{6.6}{34}}
\expandafter\def\csname r@eq:dimer-acrit-epsN\endcsname{{6.7}{34}}
\expandafter\def\csname r@sec:nbd\endcsname{{6.4}{35}}
\expandafter\def\csname r@sec:gamma\endcsname{{6.5}{36}}
\expandafter\def\csname r@fig:Gamma'\endcsname{{6.2}{36}}
\expandafter\def\csname r@sec:curve\endcsname{{6.6}{36}}
\expandafter\def\csname r@sec:pseudo\endcsname{{7}{38}}
\expandafter\def\csname r@lemma:pd-m2\endcsname{{7.1}{39}}
\expandafter\def\csname r@lemma:m1m2\endcsname{{7.2}{39}}
\expandafter\def\csname r@lemma:pd-obtuse\endcsname{{7.4}{39}}
\expandafter\def\csname r@lemma:pd-n2\endcsname{{7.5}{40}}
\expandafter\def\csname r@lemma:pd-1\endcsname{{7.6}{40}}
\expandafter\def\csname r@lemma:pd-b\endcsname{{7.7}{40}}
\expandafter\def\csname r@sec:dimer-main\endcsname{{8}{41}}
\expandafter\def\csname r@thm:dimer\endcsname{{8.1}{41}}
\expandafter\def\csname r@lemma:singleton\endcsname{{8.2}{41}}
\expandafter\def\csname r@lemma:pd-nra-dimer\endcsname{{8.3}{41}}
\expandafter\def\csname r@lemma:m2-dimer\endcsname{{8.4}{41}}
\expandafter\def\csname r@lemma:m1-dimer\endcsname{{8.5}{42}}
\expandafter\def\csname r@lemma:t2b-nonobtuse\endcsname{{8.6}{42}}
\expandafter\def\csname r@lemma:n2b\endcsname{{8.7}{42}}
\expandafter\def\csname r@lemma:m2b\endcsname{{8.8}{43}}
\expandafter\def\csname r@eqn:2-ac\endcsname{{8.9}{43}}
\expandafter\def\csname r@lemma:rab-sequence\endcsname{{8.10}{43}}
\expandafter\def\csname r@lemma:notlong\endcsname{{8.11}{44}}
\expandafter\def\csname r@lemma:3m2\endcsname{{8.12}{44}}
\expandafter\def\csname r@lemma:Nb\endcsname{{8.13}{44}}
\expandafter\def\csname r@lemma:no-ao\endcsname{{8.14}{45}}
\expandafter\def\csname r@sec:obtuse\endcsname{{9}{45}}
\expandafter\def\csname r@rem:delaunay\endcsname{{9.1}{45}}
\expandafter\def\csname r@lemma:areta\endcsname{{9.2}{45}}
\expandafter\def\csname r@fig:areta\endcsname{{9.1}{46}}
\expandafter\def\csname r@lemma:obtuse-area-monotone\endcsname{{9.3}{46}}
\expandafter\def\csname r@lemma:t2b\endcsname{{9.4}{46}}
\expandafter\def\csname r@lemma:sequence\endcsname{{9.5}{47}}
\expandafter\def\csname r@lemma:c-weak\endcsname{{9.6}{47}}
\expandafter\def\csname r@eqn:c-weak\endcsname{{9.7}{47}}
\expandafter\def\csname r@lemma:card4\endcsname{{9.8}{47}}
\expandafter\def\csname r@thm:obtuse\endcsname{{9.9}{47}}
\expandafter\def\csname r@lemma:obtuse-source\endcsname{{9.10}{48}}
\expandafter\def\csname r@lemma:tb\endcsname{{9.11}{48}}
\expandafter\def\csname r@lemma:nonobtuse-radius-area\endcsname{{9.12}{49}}
\expandafter\def\csname r@eqn:n\endcsname{{9.13}{50}}
\expandafter\def\csname r@eqn:173\endcsname{{9.14}{54}}
\expandafter\def\csname r@eqn:968\endcsname{{9.15}{54}}
\expandafter\def\csname r@sec:nonobtuse\endcsname{{10}{55}}
\expandafter\def\csname r@lemma:m2\endcsname{{10.1}{56}}
\expandafter\def\csname r@lemma:2small\endcsname{{10.2}{56}}
\expandafter\def\csname r@fig:not-large\endcsname{{10.1}{57}}
\expandafter\def\csname r@lemma:n2-0\endcsname{{10.3}{57}}
\expandafter\def\csname r@lemma:T1-all\endcsname{{10.4}{57}}
\expandafter\def\csname r@lemma:4\endcsname{{10.5}{57}}
\expandafter\def\csname r@lemma:3\endcsname{{10.6}{58}}
\expandafter\def\csname r@lemma:2\endcsname{{10.7}{58}}
\expandafter\def\csname r@eqn:2\endcsname{{10.8}{59}}
\expandafter\def\csname r@sec:coords\endcsname{{A}{60}}
\expandafter\def\csname r@sec:appendix\endcsname{{A}{60}}
\expandafter\def\csname r@sec:two\endcsname{{A.1}{60}}
\expandafter\def\csname r@fig:theta\endcsname{{A.1}{60}}
\expandafter\def\csname r@eqn:sym\endcsname{{A.1}{60}}
\expandafter\def\csname r@sec:xalpha\endcsname{{A.2}{60}}
\expandafter\def\csname r@fig:extended\endcsname{{A.2}{62}}
\expandafter\def\csname r@sec:banana\endcsname{{A.3}{62}}
\expandafter\def\csname r@fig:banana2\endcsname{{A.3}{62}}
\expandafter\def\csname r@sec:pet\endcsname{{A.4}{63}}
\expandafter\def\csname r@fig:p-triangle\endcsname{{A.4}{64}}
\expandafter\def\csname r@sec:shared\endcsname{{A.7}{64}}
\expandafter\def\csname r@eqn:abc\endcsname{{A.4}{66}}
\expandafter\def\csname r@fig:cord-delta\endcsname{{A.5}{66}}
\expandafter\def\csname r@sec:pinwheel\endcsname{{A.9}{67}}
\expandafter\def\csname r@fig:cord-pinwheel\endcsname{{A.6}{67}}
\expandafter\def\csname r@fig:cord-L\endcsname{{A.7}{68}}
\expandafter\def\csname r@sec:one-ljx\endcsname{{A.10.3}{68}}
\expandafter\def\csname r@fig:cord-T\endcsname{{A.8}{69}}
\expandafter\def\csname r@sec:one-tjx\endcsname{{A.11.1}{69}}
\expandafter\def\csname r@sec:pint\endcsname{{A.12}{70}}
\expandafter\def\csname r@fig:cord-pint\endcsname{{A.9}{70}}
\expandafter\def\csname r@lemma:0605\endcsname{{A.5}{70}}
\expandafter\def\csname r@sec:one-pintx\endcsname{{A.12.1}{71}}
\expandafter\def\csname r@appendix:cc\endcsname{{B}{71}}
\expandafter\def\csname r@sec:auto-diff\endcsname{{B.2}{72}}
\expandafter\def\csname r@sec:mitm\endcsname{{B.3}{73}}
\expandafter\def\csname r@eqn:assembly\endcsname{{B.1}{74}}
\expandafter\def\csname r@sec:calc-prep\endcsname{{B.4}{74}}
\expandafter\def\csname r@lemma:peripheral-o2c\endcsname{{B.2}{75}}
\expandafter\def\csname r@sec:pseudo1\endcsname{{B.4.1}{75}}
\expandafter\def\csname r@sec:pseudo2\endcsname{{B.4.2}{76}}
\expandafter\def\csname r@eqn:pseudo2-constraint\endcsname{{B.3}{76}}
\expandafter\def\csname r@sec:pseudo-area\endcsname{{B.4.3}{77}}
\expandafter\def\csname r@sec:pseudo-area3\endcsname{{B.4.4}{77}}
\expandafter\def\csname r@sec:dimer-isosc\endcsname{{B.4.5}{78}}
\expandafter\def\csname r@sec:large\endcsname{{B.4.6}{78}}
\expandafter\def\csname r@sec:pent3\endcsname{{B.4.7}{79}}
\expandafter\def\csname r@sec:pent4\endcsname{{B.4.8}{80}}
\makeatother
\renewcommand{\thefigure}{\currentfigurenumber}
\def\tikzfig#1#2#3{\begin{figure}[p]\centering\begin{tikzpicture}#3\end{tikzpicture}\caption{#2}\end{figure}}
\begin{center}\Large Pentagon figure preservation gallery\end{center}
This document renders the original TikZ figure calls without altering their
geometries or rounded display coordinates. Numbers are the compiled paper
numbers. These are illustrations, not rigorous numerical input. Historical
hypothesis-search scratchpads are not executed.
\clearpage
\def\currentfigurenumber{1.1}
% Exact TikZ call from pentagon_v005.tex; no geometric modification.
\tikzfig{double-lattice}{The pentagonal ice-ray.  All figures show
pentagons in red and Delaunay triangles in blue.}
{
\pent{0.0}{0.0}{-90}{0.4};  %
\pent{0.0}{0.724}{-90}{0.4};  %
\pent{0.0}{2*0.724}{-90}{0.4};  
\pent{0.57}{-0.3854}{90}{0.4};  %
\pent{0.57}{-0.3854+0.724}{90}{0.4}; 
\pent{0.57}{-0.3854+2*0.724}{90}{0.4};  
\pent{0.57}{-0.3854+3*0.724}{90}{0.4};  
\pent{-0.57}{-0.3854}{90}{0.4};  %
\pent{-0.57}{-0.3854+0.724}{90}{0.4}; 
\pent{-0.57}{-0.3854+2*0.724}{90}{0.4};  
\pent{-0.57}{-0.3854+3*0.724}{90}{0.4};  
\pent{2*0.57}{0.0}{-90}{0.4};  %
\pent{2*0.57}{0.724}{-90}{0.4};  %
\pent{2*0.57}{2*0.724}{-90}{0.4}; 
\pent{3*0.57}{-0.3854}{90}{0.4};  %
\pent{3*0.57}{-0.3854+0.724}{90}{0.4}; 
\pent{3*0.57}{-0.3854+2*0.724}{90}{0.4};  
\pent{3*0.57}{-0.3854+3*0.724}{90}{0.4};   
}

\clearpage
\def\currentfigurenumber{2.2}
% Exact TikZ call from pentagon_v005.tex; no geometric modification.
\tikzfig{KKtriangle}{Ice-ray triangles and ice-ray dimers.
  All Delaunay triangles in the pentagonal ice-ray are congruent and
  have area $\acrit$.  We call them ice-ray triangles. 
  The pentagonal ice-ray can be partitioned into pairs of Delaunay
  triangles (called ice-ray dimers) in which the two triangles in each
  dimer share their common longest edge.  
}
{
\begin{scope}
\pentpink{0.0}{0.0}{-90}{0.4}{red!30};  %
\pentpink{0.0}{0.724}{-90}{0.4}{red!30};  %
\pentpink{0.0}{2*0.724}{-90}{0.4}{red!30};  
\pentpink{0.57}{-0.3854}{90}{0.4}{red!30};  %
\pentpink{0.57}{-0.3854+0.724}{90}{0.4}{red!30}; 
\pentpink{0.57}{-0.3854+2*0.724}{90}{0.4}{red!30};  
\pentpink{0.57}{-0.3854+3*0.724}{90}{0.4}{red!30};  
\pentpink{-0.57}{-0.3854}{90}{0.4}{red!30};  %
\pentpink{-0.57}{-0.3854+0.724}{90}{0.4}{red!30}; 
\pentpink{-0.57}{-0.3854+2*0.724}{90}{0.4}{red!30};  
\pentpink{-0.57}{-0.3854+3*0.724}{90}{0.4}{red!30};  
\pentpink{2*0.57}{0.0}{-90}{0.4}{red!30};  %
\pentpink{2*0.57}{0.724}{-90}{0.4}{red!30};  %
\pentpink{2*0.57}{2*0.724}{-90}{0.4}{red!30}; 
\pentpink{3*0.57}{-0.3854}{90}{0.4}{red!30};  %
\pentpink{3*0.57}{-0.3854+0.724}{90}{0.4}{red!30}; 
\pentpink{3*0.57}{-0.3854+2*0.724}{90}{0.4}{red!30};  
\pentpink{3*0.57}{-0.3854+3*0.724}{90}{0.4}{red!30};   
\coordinate (T) at (0.5706,-0.38541);
\coordinate (U) at (0,0.7236);
\coordinate (T') at ($(T)+(U)$);
\draw[blue] (0,0) -- ++ ($2*(U)$);
\draw[blue] ($-1.0*(T')$) -- ++ ($3*(U)$);
\draw[blue] (T) -- ++ ($3*(U)$);
\draw[blue] ($(T)+(T')$) -- ++ ($2*(U)$);
\draw[blue] ($(T)+(T')+(T)$) -- ++ ($3*(U)$);
\draw[blue] ($-1.0*(T)$) -- ++ ($2.0*(T)$);
\draw[blue] ($-1.0*(T) + (U)$) -- ++ ($4.0*(T)$);
\draw[blue] ($-1.0*(T) + 2.0*(U)$) -- ++ ($4.0*(T)$);
\draw[blue] ($1.0*(T) + 3.0*(U)$) -- ++ ($2.0*(T)$);
\draw[blue] ($(T)$) -- ++ ($2.0*(T')$);
\draw[blue] ($1.0*(T) +1.0*(U) -2.0 *(T')$) -- ++ ($4.0*(T')$);
\draw[blue] ($1.0*(T) +2.0*(U) -2.0 *(T')$) -- ++ ($4.0*(T')$);
\draw[blue] ($1.0*(T) +3.0*(U) -2.0 *(T')$) -- ++ ($2.0*(T')$);
\end{scope}
\begin{scope}[xshift=5cm]
\pentpink{0.0}{0.0}{-90}{0.4}{red!30};  %
\pentpink{0.0}{0.724}{-90}{0.4}{red!30};  %
\pentpink{0.0}{2*0.724}{-90}{0.4}{red!30};  
\pentpink{0.57}{-0.3854}{90}{0.4}{red!30};  %
\pentpink{0.57}{-0.3854+0.724}{90}{0.4}{red!30}; 
\pentpink{0.57}{-0.3854+2*0.724}{90}{0.4}{red!30};  
\pentpink{0.57}{-0.3854+3*0.724}{90}{0.4}{red!30};  
\pentpink{-0.57}{-0.3854}{90}{0.4}{red!30};  %
\pentpink{-0.57}{-0.3854+0.724}{90}{0.4}{red!30}; 
\pentpink{-0.57}{-0.3854+2*0.724}{90}{0.4}{red!30};  
\pentpink{-0.57}{-0.3854+3*0.724}{90}{0.4}{red!30};  
\pentpink{2*0.57}{0.0}{-90}{0.4}{red!30};  %
\pentpink{2*0.57}{0.724}{-90}{0.4}{red!30};  %
\pentpink{2*0.57}{2*0.724}{-90}{0.4}{red!30}; 
\pentpink{3*0.57}{-0.3854}{90}{0.4}{red!30};  %
\pentpink{3*0.57}{-0.3854+0.724}{90}{0.4}{red!30}; 
\pentpink{3*0.57}{-0.3854+2*0.724}{90}{0.4}{red!30};  
\pentpink{3*0.57}{-0.3854+3*0.724}{90}{0.4}{red!30};   
\coordinate (T) at (0.5706,-0.38541);
\coordinate (U) at (0,0.7236);
\coordinate (T') at ($(T)+(U)$);
\draw[blue] ($-1.0*(T)$) -- ++ ($2.0*(T)$);
\draw[blue] ($-1.0*(T) + (U)$) -- ++ ($4.0*(T)$);
\draw[blue] ($-1.0*(T) + 2.0*(U)$) -- ++ ($4.0*(T)$);
\draw[blue] ($1.0*(T) + 3.0*(U)$) -- ++ ($2.0*(T)$);
\draw[blue] ($(T)$) -- ++ ($2.0*(T')$);
\draw[blue] ($1.0*(T) +1.0*(U) -2.0 *(T')$) -- ++ ($4.0*(T')$);
\draw[blue] ($1.0*(T) +2.0*(U) -2.0 *(T')$) -- ++ ($4.0*(T')$);
\draw[blue] ($1.0*(T) +3.0*(U) -2.0 *(T')$) -- ++ ($2.0*(T')$);
\draw[blue] ($-1.0*(T) + 2.0*(U)$) -- ++ ($-1.0*(T')$) -- ++ (T)
  -- ++ ($-1.0*(T')$) -- ++ (T)
  -- ++ ($-1.0*(T')$) -- ++ (T);
\draw[blue] ($3.0*(T) + 4.0*(U)$) -- ++ (T) -- ++ ($-1.0*(T')$) 
  -- ++ (T)
  -- ++ ($-1.0*(T')$) 
  -- ++ (T)
  -- ++ ($-1.0*(T')$);
\end{scope}
}

\clearpage
\def\currentfigurenumber{2.3}
% Exact TikZ call from pentagon_v005.tex; no geometric modification.
\tikzfig{subcritical}{Subcritical triangles.
Experimentally, the first triangle has the smallest area among all
acute Delaunay triangles.  The second triangle is a cloverleaf that
has two edges of minimal length $2\kappa$.  Experimentally, the
third triangle minimizes the longest edge among subcritical acute Delaunay triangles.
It is equilateral with edge length about $1.72256$.}
{
\begin{scope}[scale=0.5]
\threepent{0.00}{0.00}{85.23}{0.93}{1.33}{265.23}{1.86}{0.00}{229.23};
\draw (0.5,-1.5) node {$\area\approx 1.23719$};
\end{scope}
\begin{scope}[scale=0.5,xshift=5cm]
\threepent{0.00}{0.00}{90.00}{0.95}{1.31}{270.00}{1.90}{0.00}{234.00};
\draw (0.5,-1.5) node {$\area\approx 1.24$};
\end{scope}
\begin{scope}[scale=0.5,xshift=10cm]
\threepent{0.00}{0.00}{96.03}{0.86}{1.49}{264.03}{1.72}{0.00}{216.03};
\draw (0.5,-1.5) node {$\area\approx 1.285$};
\end{scope}
\begin{scope}[scale=0.5,xshift=15cm] %
\threepent{0.00}{0.00}{109.36}{0.83}{1.59}{235.36}{1.62}{0.00}{145.36};
\draw (0.5,-1.5) node {$\area\approx 1.286$};
\end{scope}
}

\clearpage
\def\currentfigurenumber{2.4}
% Exact TikZ call from pentagon_v005.tex; no geometric modification.
\tikzfig{KK2min}{Ice-ray triangle deformation.
The ice-ray triangle is not a local minimum of the area function.
The  ice-ray triangle (left) can be continuously deformed
along an area decreasing curve to the subcritical acute triangle of 
numerically minimum area (right). }
{
\begin{scope}[scale=0.5]
\threepent{0.00}{0.00}{72.00}{0.96}{1.43}{252.00}{1.81}{0.00}{216.00};
\draw (0.5,-1.5) node {$\area=\acrit$};
\end{scope}
\begin{scope}[scale=0.5,xshift=6cm]
\threepent{0.00}{0.00}{79.49}{0.93}{1.37}{259.49}{1.82}{0.00}{223.49};
\draw (0.5,-1.5) node {$\area\approx 1.248$};
\end{scope}
\begin{scope}[scale=0.5,xshift=12cm]
\threepent{0.00}{0.00}{85.23}{0.93}{1.33}{265.23}{1.86}{0.00}{229.23};
\draw (0.5,-1.5) node {$\area\approx 1.23719$};
\end{scope}
}

\clearpage
\def\currentfigurenumber{2.5}
% Exact TikZ call from pentagon_v005.tex; no geometric modification.
\tikzfig{Gamma}{Ice-ray dimer deformation.  The ice-ray
  dimer (center) admits a shear motion that preserves all edges of
  contact.  This deformation increases the area of the dimer.  It is
  obvious by the symmetry of the figures on the left and right that
  the area function along this deformation has a critical point at the
  ice-ray dimer.  It is known that the ice-ray dimer is a local
  minimum of the area function \cite{Kus}.  }
{
\begin{scope}[scale=0.5]
\threepent{0.00}{0.00}{72.00}{0.96}{1.43}{252.00}{1.81}{0.00}{216.00};
\threepent{0.00}{0.00}{-72.00}{0.96}{-1.43}{-252.00}{1.81}{0.00}{-216.00};
\draw (0.5,-3.5) node {$\area=2\acrit$};
\end{scope}
\begin{scope}[scale=0.5,xshift=6cm]
\threepent{0.00}{0.00}{79.49}{0.93}{1.37}{259.49}{1.82}{0.00}{223.49};
\threepent{0.00}{0.00}{-64.51}{1.03}{-1.50}{-244.51}{1.82}{0.00}{-208.51};
\draw (0.5,-3.5) node {$\area\approx 2\acrit+0.03$};
\end{scope}
\begin{scope}[scale=0.5,xshift=-6cm]
\threepent{0.00}{0.00}{-79.49}{0.93}{-1.37}{-259.49}{1.82}{0.00}{-223.49};
\threepent{0.00}{0.00}{64.51}{1.03}{1.50}{244.51}{1.82}{0.00}{208.51};
\draw (0.5,-3.5) node {$\area\approx 2\acrit+0.03$};
\end{scope}
}

\clearpage
\def\currentfigurenumber{2.6}
% Exact TikZ call from pentagon_v005.tex; no geometric modification.
\tikzfig{pseudo-dimer}{Pseudo-dimer.  There exist pairs of acute
  triangles $(T_1,T_0)$ such that the sum of their two areas is at
  most $2\acrit$ and such that (1) $T_1$ is subcritical and is
  adjacent to $T_0$ along the longest edge of $T_1$, but such that (2)
  the longest edge of $T_0$ is not the edge shared with $T_1$.  We
  call such pairs {\it pseudo-dimers}.    The illustrated
  pseudo-dimer has area about $2\acrit - 10^{-5}$, and the triangle $T_0$
  has longest edge length about $1.84$.  The modification
  term $\epsM = 0.008$ is based on this and closely related examples
  of pseudo-dimers.  Pseudo-dimers add
  significant complications to the proof. }
{
\begin{scope}[scale=1.2,rotate=90]
\threepent{0.00}{0.00}{93.42}{0.90}{1.48}{259.93}{1.73}{0.00}{214.93};
\threepent{0.00}{0.00}{-50.58}{1.10}{-1.49}{-181.07}{1.73}{0.00}{-217.07};
\draw (1.3,-0.1) node {$T_1\Rightarrow T_0$};
\draw (0.8,-1) node {$\Searrow$};
\end{scope}
}

\clearpage
\def\currentfigurenumber{2.7}
% Exact TikZ call from pentagon_v005.tex; no geometric modification.
\tikzfig{obtuse}{Obtuse Delaunay triangles can have small area.  A
  triangle with circumradius about $2$ and area about $0.98$
  is shown. The adjacent Delaunay triangle cannot have a vertex
  inside the circumcircle of the first.  The neighbor of a subcritical
  obtuse Delaunay triangle tends to have large area.}
{
\begin{scope}[scale=0.5]
\threepent{0}{0}{90}
{1.486}{-0.661}{-90}
{-1.486}{-0.661}{-90};
\end{scope}
\begin{scope}[scale=0.5,xshift=6cm]
\draw[blue] (0,0)-- (1.486,-0.661) -- (-1.486,-0.661) -- cycle;
\draw[blue] (1.486,-0.661) -- ($(0,-2) + (-6:2)$) -- (-1.486,-0.661);
\draw (0,-2) circle (2cm);
\end{scope}
}

\clearpage
\def\currentfigurenumber{2.8}
% Exact TikZ call from pentagon_v005.tex; no geometric modification.
\tikzfig{visual}{The area modification $b(T)$ can
be visualized as a small semicircular 
knob (of area $\epsN$) or triangular notch (of area $\epsM$) 
that transfers area
from one Delaunay triangle to an adjacent one. 
}
{
\begin{scope}[scale=2,xshift=0cm]
\coordinate (T) at (0.5706,-0.38541);
\coordinate (U) at (0,0.7236);
\coordinate (T') at ($(T)+(U)$);
\draw[blue] (U) -- ++ (T) -- ++ ($-1.0*(T')$) -- ++ ($-1.0*(T)$) -- ++ (T');
\draw[blue] (0,0) -- ++ ($0.5*(U)-(0,0.08)$);
\draw[blue] (0, 0.5*0.7236 - 0.08)
  arc[start angle=-90, end angle=90, radius=0.08];
\draw[blue] ($0.5*(U)+(0,0.08)$) -- (U);
\end{scope}
\begin{scope}[scale=2,xshift=1.5cm]
\coordinate (T) at (0.5706,-0.38541);
\coordinate (U) at (0,0.7236);
\coordinate (T') at ($(T)+(U)$);
\draw[blue] (U) -- ++ (T) -- ++ ($-1.0*(T')$) -- ++ ($-1.0*(T)$) -- ++ (T');
\draw[blue] (0,0) -- ++ ($0.5*(U)-(0,0.055/2.0)$) -- ++ (30:0.055) --++ (150:0.055)--(U) ;
\end{scope}
}

\clearpage
\def\currentfigurenumber{3.1}
% Exact TikZ call from pentagon_v005.tex; no geometric modification.
\tikzfig{receptor}{Pointer and receptor pairs of pentagons.  
In each pair, the pentagon on the
left can be considered a pointer pentagon, 
with receptor on the right.  The first pair is nondegenerate,
and the other two pairs are degenerate.}{
[scale=0.5]
\pent{0}{0}{0}{1};
\pent{1.8}{0}{0}{1};
\pent{5}{0}{0}{1};
\pent{7}{0}{180}{1};
\pent{10+0.176}{0.243}{0}{1};  %
\pent{12-0.172}{-0.243}{180}{1};
}

\clearpage
\def\currentfigurenumber{3.2}
% Exact TikZ call from pentagon_v005.tex; no geometric modification.
\tikzfig{3C-type}{Types of $3C$-contact from top-to-bottom:
a pinwheel, a pin-$T$ junction, 
a $\Delta$-junction, an $LJ$-junction, and a $TJ$-junction.
In the first two types, arrows form a directed cycle. In the other
three types, both arrows point outward from the pentagon on top.
The second column gives \emph{distorted} representations of the
inner polygons.
In the pin-$T$, the inner polygon, which is
a nonconvex pentagon, is difficult to discern because of a 
necessarily small angle and
short edge.   
}%
{
[scale=0.8] 
\begin{scope}[yshift=20cm]
\threepentnoD{0.00}{0.00}{46.69}{0.82}{1.53}{218.09}{1.73}{0.00}{163.28};
\coordinate (A) at (0,0);
\coordinate (B) at (0.82,1.53);
\coordinate (C) at (1.73,0);
\draw[->] ($0.8*(A) + 0.2*(B)$) -- ($0.2*(A) + 0.8*(B)$);
\draw[->] ($0.8*(B) + 0.2*(C)$) -- ($0.2*(B) + 0.8*(C)$);
\draw[->] ($0.8*(C) + 0.2*(A)$) -- ($0.2*(C) + 0.8*(A)$);
\end{scope}
\begin{scope}[xshift=4.5cm,yshift=20cm,scale=0.8]
\draw[red] (0,0) -- ++ (30:1cm)  
-- ++ (150:1cm) node[black,anchor=south] {\hbox{pinwheel}} -- (0,0);
\end{scope}
\begin{scope}[yshift=17.5cm]
\threepentnoD{0.00}{0.00}{90}{1.40}{-1.377}{0}{-0.35}{-1.628}{-18.89};
\coordinate (A) at (0,0);
\coordinate (B) at (1.40,-1.377);
\coordinate (C) at (-0.35,-1.628);
\draw[<-] ($0.8*(A) + 0.2*(B)$) -- ($0.2*(A) + 0.8*(B)$);
\draw[<-] ($0.8*(B) + 0.2*(C)$) -- ($0.2*(B) + 0.8*(C)$);
\draw[<-] ($0.8*(C) + 0.2*(A)$) -- ($0.2*(C) + 0.8*(A)$);
\end{scope}
\begin{scope}[xshift=5.5cm,yshift=15.5cm,scale=0.8]
\draw[red] (0,0) -- ++ (90:2cm) 
-- ++ (245:0.5cm) 
-- node[anchor=south,black] {pin-$T$} ++ (180:1.5cm) -- ++ (350:1.2cm) -- (0,0);
\end{scope}
\begin{scope}[yshift=12cm]
\threepentnoD{0.00}{0.00}{-5.16}{0.99}{1.70}{235.38}{1.98}{0.00}{183.43};
\coordinate (A) at (0,0);
\coordinate (B) at (0.99,1.7);
\coordinate (C) at (1.98,0);
\draw[<-] ($0.8*(A) + 0.2*(B)$) -- ($0.2*(A) + 0.8*(B)$);
\draw[->] ($0.8*(B) + 0.2*(C)$) -- ($0.2*(B) + 0.8*(C)$);
\draw[->] ($0.8*(C) + 0.2*(A)$) -- ($0.2*(C) + 0.8*(A)$);
\end{scope}
\begin{scope}[xshift=5cm,yshift=12cm,scale=1.3]
\draw[red] (0,0) -- ++ (60:1cm)  
-- node[black,anchor=south] {$\Delta$} ++ (180:1cm)  -- (0,0);
\end{scope}
\begin{scope}[yshift=8cm]
\threepentnoD{0.00}{0.00}{80.86}{0.97}{1.58}{232.21}{1.82}{0.00}{223.24};
\coordinate (A) at (0,0);
\coordinate (B) at (0.97,1.58);
\coordinate (C) at (1.82,0);
\draw[<-] ($0.8*(A) + 0.2*(B)$) -- ($0.2*(A) + 0.8*(B)$);
\draw[->] ($0.8*(B) + 0.2*(C)$) -- ($0.2*(B) + 0.8*(C)$);
\draw[<-] ($0.8*(C) + 0.2*(A)$) -- ($0.2*(C) + 0.8*(A)$);
\end{scope}
\begin{scope}[xshift=5cm,yshift=8cm,scale=0.7]
\draw[red] (0,0) -- ++ (95:0.8cm)  
-- ++ (30:1cm) 
-- node[black,anchor=south] {\hbox{LJ}} ++ (180:2cm) 
-- (0,0);
\end{scope}
\begin{scope}[yshift=4cm]
\threepentnoD{0.00}{0.00}{114.48}{0.90}{1.59}{237.18}{1.66}{0.00}{219.24};
\coordinate (A) at (0,0);
\coordinate (B) at (0.9,1.59);
\coordinate (C) at (1.66,0);
\draw[<-] ($0.8*(A) + 0.2*(B)$) -- ($0.2*(A) + 0.8*(B)$);
\draw[->] ($0.8*(B) + 0.2*(C)$) -- ($0.2*(B) + 0.8*(C)$);
\draw[<-] ($0.8*(C) + 0.2*(A)$) -- ($0.2*(C) + 0.8*(A)$);
\end{scope}
\begin{scope}[xshift=5cm,yshift=3.5cm,scale=0.65]
\draw[red] (0,0) -- ++ (90:1.8cm)  
-- ++ (20:1.2cm) 
-- node[black,anchor=south] {\hbox{TJ}} ++ (180:2.5cm) 
-- ++ (345:1.2cm)
-- (0,0);
\end{scope}
}

\clearpage
\def\currentfigurenumber{3.3}
% Exact TikZ call from pentagon_v005.tex; no geometric modification.
\tikzfig{clover}{A cloverleaf (degenerate pinwheel).  The Delaunay triangle
in this particular example is not subcritical.}
{
[scale=0.6]
\begin{scope}[xshift=4cm]
\threepentnoD{0.00}{0.00}{31.70}{0.76}{1.52}{203.11}{1.70}{0.00}{148.30};
\end{scope}
}

\clearpage
\def\currentfigurenumber{3.4}
% Exact TikZ call from pentagon_v005.tex; no geometric modification.
\tikzfig{sep}{A line through the center 
of the middle pentagon $B$
through one of its vertices 
separates the two extremal 
pentagons $A$ and $C$.}{
[scale=0.5]
\pen{0}{0}{90};
\pen{1.72}{0.56}{90};  %
\pen{-1.72}{0.56}{90};  
\draw (0,-1.5) -- (0,1.5) node[black,anchor=west] {$\v$};
\begin{scope}[xshift=6cm]
\pen{0}{0}{90};
\pen{1.72}{0.56}{90};  
\pen{-1.72}{-0.56}{90};  
\draw (0,-1.5) -- (0,1.5) node[black,anchor=west] {$\v$};
\end{scope}
\begin{scope}[xshift=12cm]
\pen{0}{0}{90};
\pen{1.72}{-0.56}{90};  
\pen{-1.72}{-0.56}{90};  
\draw (0,-1.5) -- (0,1.5) node[black,anchor=west] {$\v$};
\end{scope}
}

\clearpage
\def\currentfigurenumber{3.5}
% Exact TikZ call from pentagon_v005.tex; no geometric modification.
\tikzfig{delta}{In type $\Delta$, a line separates pentagon 
$B$ from
the other two pentagons. 
The second figure (which is degenerate of type $LJ$) does not have
type $\Delta$.}{
[scale=0.5]
\pen{0}{0}{90};
\draw (0,0) node {$B$};
\draw (-2.0,-0.809) node[above] {$\lambda$} --  (2.0,-0.809);
\pen{-1.05}{-2*0.809}{-90};
\pen{1.9 - 1.05}{-2*0.809}{-90};
\begin{scope}[xshift=6cm]
\pen{0}{0}{90};
\draw (0,0) node {$B$};
\draw (-2.0,-0.809) node[above] {$\lambda$} --  (2.0,-0.809);
\pen{-0.428}{-2*0.809}{-90};
\pen{1.36}{ -1.44}{ 212.704};
\end{scope}
}

\clearpage
\def\currentfigurenumber{3.6}
% Exact TikZ call from pentagon_v005.tex; no geometric modification.
\tikzfig{pin1}{A distorted pin-$1$ configuration.}{
\begin{scope}[scale=1.0]
\pen{0}{0}{-90};
\pen{-0.974}{ 1.543}{ 32.70422};
\draw[red] ++ (-0.974,1.543) 
  ++ (0.842,0.532) node[black,anchor=west] {$\v_{BC}$} -- 
  ++ (115:0.2) -- ++ (-65:1.6) node[black,anchor=south west] {$\v_{CA}$}
   -- ++ (7:1.17);
\draw (0,0) node {$A$};
\draw (0,0) + (3*72-90:1) node[anchor=north west] {$\v_{AB}$};
\draw (-0.974,1.543) node {$B$};
\draw (1.1,1.543) node {$C$};
\smalldot{3*72-90:1}; %
\draw (-0.974,1.543)+(32.7-72:1) node[anchor=south west]{$\v$};
\smalldot{$(-0.974,1.543)+(32.7-72:1)$}; %
\smalldot{$(-0.974,1.543)+(32.7:1)$}; %
\smalldot{$(-0.974,1.543)+(0.842,0.532)+(115:0.2)+(-65:1.6)$}; %
\end{scope}
}

\clearpage
\def\currentfigurenumber{5.1}
% Exact TikZ call from pentagon_v005.tex; no geometric modification.
\tikzfig{large}{A Delaunay triangle $T$ 
with a large angle $\alpha$ along $(T,e)$.}
{
\begin{scope}[scale=0.8]
\threepent{0.00}{0.00}{50.58}{1.10}{1.49}{181.07}{1.73}{0.00}{217.07};
\draw (0,0) node[anchor=north] {$A$};
\draw (0,0) node[anchor=south] {$e\ $};
\draw (1.1,1.49) node[anchor=south] {$B$};
\draw[black] (50.58:1) -- ++ ($2*(50.58+72:1) - 2*(50.58:1)$);
\draw[black] (50.58:1) -- ++ ($2*(1.1,1.49)+2*(181.07:1) - 2*(50.58:1)$);
\draw (-0.8,1.4) node {$\alpha$};
\end{scope}
}

\clearpage
\def\currentfigurenumber{6.1}
% Exact TikZ call from pentagon_v005.tex; no geometric modification.
\tikzfig{slider}{The center $\c_B$ of the nonobtuse triangle lands in
  the first quadrant and gives a triangle $(\c_A,\c_B,\c_C)$ of height
  and base both at least $2\kappa$.}
{
[scale=1.0]
\draw (0,0) -- (2,0);
\draw (0,0) -- (0,2);
\coordinate (V) at ($(0,0) + (-90+2*72+10:1)   $);
\draw (V) -- ++ (10:1.5) node[anchor=west] {$\lambda$};
\draw (V) -- ++ (180+10:1.5);
\coordinate (B) at ($(V) + (10:0.5) + (90+10- 2*72 - 180:1)$);
\draw (B) -- ++ (10:1.5) node[anchor=west] {$\lambda'$};
\draw (B) -- ++ (180+10:1.5);
\penp{B}{90+10};
\pen{0}{0}{-90+10};
\smalldot{B};
\smalldot{0,0};
\smalldot{2,0};
\draw (B) node [anchor=south west] {$\c_B$};
\draw (0,0) node [anchor=north] {$\c_A$};
\draw (2,0) node [anchor=north] {$\c_C$};
}

\clearpage
\def\currentfigurenumber{6.2}
% Exact TikZ call from pentagon_v005.tex; no geometric modification.
\tikzfig{Gamma'}{Three configurations along the curve $\Gamma$.}
{
\begin{scope}[scale=0.5,xshift=-6cm]
\threepentnoD{0.00}{0.00}{79.87}{0.93}{1.36}{259.87}{1.83}{0.00}{223.87};
\threepentnoD{0.00}{0.00}{-64.13}{1.03}{-1.50}{-244.13}
{1.83}{0.00}{-208.13};
\draw (0,0) node {$A$};
\draw (1.83,0) node {$C$};
\draw (0.5,-3.5) node {$t= 0.25$};
\smalldot{79.87-72:1};
\smalldot{$(1.83,0)+(223.87-3*72+180:0.809)$};
\end{scope}
\begin{scope}[scale=0.5]
\threepentnoD{0.00}{0.00}{72.00}{0.96}{1.43}{252.00}{1.81}{0.00}{216.00};
\threepentnoD{0.00}{0.00}{-72.00}{0.96}{-1.43}{-252.00}{1.81}{0.00}{-216.00};
\draw (0,0) node {$A$};
\draw (1.81,0) node {$C$};
\smalldot{1,0};
\draw (0.5,-3.5) node {$t=0$};
\end{scope}
\begin{scope}[scale=0.5,xshift=6cm]
\threepentnoD{0.00}{0.00}{-79.87}{0.93}{-1.36}{-259.87}{1.83}{0.00}{-223.87};
\threepentnoD{0.00}{0.00}{64.13}{1.03}{1.50}{244.13}{1.83}{0.00}{208.13};
\draw (0,0) node {$A$};
\draw (1.83,0) node {$C$};
\draw (0.5,-3.5) node {$t= -0.25$};
\smalldot{-79.87+72:1};
\smalldot{$(1.82,0)+(-223.87+3*72+180:0.809)$};
\end{scope}
}

\clearpage
\def\currentfigurenumber{9.1}
% Exact TikZ call from pentagon_v005.tex; no geometric modification.
\tikzfig{areta}{
There can be two positions $\r,\r'$ 
on the circumcircle for
the third vertex of the triangle.  
Here, $d_1 = \norm{\p}{\q}$.}{
[scale=1.3]
\draw (0,0) circle (1cm);
\draw (1,0) node[anchor=west] {$\p$} --  (70:1cm) node[anchor=south] {$\q$};
\draw (1,0) -- node[above=1pt] {$d_2$} (130:1cm) node[anchor=south] {$\r'$};
\draw (1,0) -- node[below=1pt] {$d_2$} (-130:1cm) node[anchor=north] {$\r$};
\draw (130:1cm) -- (-130:1cm);
\draw[thin,gray] (1,0) -- (-1,0);
\draw (100:1cm) node[above=1pt] {$\s$};
}

\clearpage
\def\currentfigurenumber{10.1}
% Exact TikZ call from pentagon_v005.tex; no geometric modification.
\tikzfig{not-large}{The angle between the nearly vertical edges of
  $A$ and $B$ is $\gamma$.  On the left, $\gamma= 0.02$ and $
  2\kappa<d_{AB}\approx 1.6296 < 1.63$.  There is almost no play in
  the configuration.  The  figure on the right does not satisfy the
  constraints of the proof of Lemma \rif{lemma:2small}, but
  illustrates the notation.}
{
\begin{scope}[scale=1]
\pen{0}{0}{35.87};
\pen{1.623}{0}{-34.72+180};
\draw (0,0) node {$A$};
\draw (1.623,0) node {$B$};
\draw (0.8,1) node {$\gamma'$};
\draw ($(1.623,0)+(-34.72+180:1)$) node[anchor=north west] {\ $\v_B$};
\smalldot{$(1.623,0)+(-34.72+180:1)$};
\end{scope}
\begin{scope}[scale=1,xshift=5cm]
\pen{0}{0}{34.86};
\pen{1.756}{0}{-17.68+180};
\draw (0,0) node {$A$};
\draw (1.756,0) node {$B$};
\draw ($(1.756,0) + (-17.68+180:1)$) --  ++ 
($1.5*(-17.68+180+72:1) - 1.5*(-17.68+180:1)$) node[anchor=east] {$\gamma$};
\draw ($(1.756,0) + (-17.68+180:1)$) --  ++ 
($1.5*(34.86-72:1) - 1.5*(34.86:1)$) ;
\coordinate (A) at (34.86:1);
\coordinate (Aa) at (34.86+72:1);
\draw ($1.21*(A) - 0.21*(Aa)$) node[anchor=south] {$\gamma'$} 
-- ($-0.4*(A) + 1.4*(Aa)$);
\draw ($(1.756,0)+(-17.68+180:1)$) node[anchor=north west] {\ $\v_B$};
\smalldot{$(1.756,0)+(-17.68+180:1)$};
\end{scope}
}

\clearpage
\def\currentfigurenumber{A.1}
% Exact TikZ call from pentagon_v005.tex; no geometric modification.
\tikzfig{theta}{coordinates for a pair of pentagons in contact}{
[scale=1.2]
\pen{0}{0}{0};
\pen{1.616}{0.598}{168.54};
\draw[blue] (0,0) node[anchor=south,black] {$B$}
 -- (1.616,0.598) node[anchor=south,black] {$A$};
\draw (0,0) node[anchor=north,black] {$\c_B$};
\smalldot{0,0};
\draw (1.616,0.598) node[anchor=north,black] {$\c_A$};
\smalldot{1.616,0.598};
\draw (72:1) node[anchor=south] {$\w_B$};
\draw (0:1) node[anchor=west] {$\v_B$};
\draw (-72:1) node[anchor=north] {$\u_B$};
\smalldot{72:1};
\smalldot{0:1};
\smalldot{-72:1};
\smalldot{0.636,0.797};
\draw (0.636,0.797) node[anchor=west] {$\w_A$}
node[anchor=north west] {$x$};
\smalldot{1.124,-0.2727};
\draw (1.124,-0.2727) node[anchor=north west] {$\u_A$}
node[anchor=north east] {$\alpha$};
}

\clearpage
\def\currentfigurenumber{A.2}
% Exact TikZ call from pentagon_v005.tex; no geometric modification.
\tikzfig{extended}{Extended coordinates give a continuous transition
  from pairs $(A,B)$ with pointer $B\to A$ to pointer $A\to B$.  The
  functions $\theta,\theta'$ are continuous in $\alpha,x_\alpha$.  The
  function $\theta$ is the angle $\angle(\c_A,\c_B,\v)$, and $\theta'
  = \angle(\c_B,\c_A,\w)$. The vertices $\v$ and $\w$ transition into
  and away from the pointer vertex.  In these figures, $\alpha$ varies
  and $x_\alpha=0.8=\norm{\v}{\w}$ is fixed.}
{
\begin{scope}[scale=1]
\pen{0}{0}{41.16};
\pen{1.745}{0}{180-16.51};
\draw[blue] (0,0) node[black,anchor=east] {$A$} 
-- (1.745,0) node[black,anchor=west] {$B$};
\draw (0.8,-2) node {$\alpha = 2\pi/5 - 0.25$};
\smalldot{41.16-72:1};
\draw (41.16-72:1) node[anchor=east] {$\v~$~};
\smalldot{$(1.745,0)+(180-16.51:1)$};
\draw ($(1.745,0)+(180-16.51:1)$)  node[anchor=west] {$\w$};
\end{scope}
\begin{scope}[scale=1,xshift=4cm]
\pen{0}{0}{49.067};
\pen{1.66}{0}{180-22.93};
\draw (0.8,-2) node {$\alpha = 2\pi/5$};
\draw[blue] (0,0) node[black,anchor=east] {$A$} 
-- (1.66,0) node[black,anchor=west] {$B$};
\smalldot{49.067-72:1};
\draw (49.067-72:1) node[anchor= east] {$\v$~};
\smalldot{$(1.66,0)+(180-22.93:1)$};
\draw ($(1.66,0)+(180-22.93:1)$)  node[anchor=west] {$~\w$};
\end{scope}
\begin{scope}[scale=1,xshift=8cm]
\pen{0}{0}{55.489};
\pen{1.745}{0}{180-30.83};
\draw (0.8,-2) node {$\alpha = 2\pi/5+0.25$};
\draw[blue] (0,0) node[black,anchor=east] {$A$} 
-- (1.745,0) node[black,anchor=west] {$B$};
\smalldot{55.489-72:1};
\draw (55.489-72:1) node[anchor=east] {$\v$};
\smalldot{$(1.745,0)+(180-30.83:1)$};
\draw ($(1.745,0)+(180-30.83:1)$)  
node[anchor=west] {$~\w$};
\end{scope}
}

\clearpage
\def\currentfigurenumber{A.3}
% Exact TikZ call from pentagon_v005.tex; no geometric modification.
\tikzfig{banana2}{With given fixed center $\c_A$ and fixed pointer
  pentagon $B$, there might be two pentagons $A$ in contact with
  $B$. They are distinguished by a boolean variable indicating whether
  the pointer $B\to A$ lies above or below the midpoint of the edge on
  $A$.}
{
\begin{scope}[scale=1,xshift=5cm]
\pen{0}{0}{29.68};
\pentpink{0}{0}{-15.887}{1}{black!50};
\pen{1.82}{0}{180-5.73};
\draw (0,0) node[black] {$A$} ;
\draw (1.82,0) node[black] {$B$};
\end{scope}
}

\clearpage
\def\currentfigurenumber{A.4}
% Exact TikZ call from pentagon_v005.tex; no geometric modification.
\tikzfig{p-triangle}{Angle conventions. 
  The edge lengths of the Delaunay triangle are
  $d_{AC}$, $d_{AB}$, and $d_{BC}$.  The angles at $\c_A$ are
  $\theta_{ABC}$, $\arc_A$, and $\theta_{ACB}$.  The angles at $\c_B$
  are $\theta_{BAC}$, $\arc_B$, and $\theta_{BCA}$. The angles at
  $\c_C$ are $\theta_{CBA}$, $\arc_C$, and $\theta_{CAB}$.}
{
\begin{scope}[scale=1.9,xshift=5cm]
\def\x{1.2}
\threepent
{0.00}{0.00}{96.88}
{\x*0.95}{\x*1.52}{252.34}
{\x*1.68}{0.00}{210.61};
\draw(0,0) node[anchor=east] {$A$};
\draw(\x*0.95,\x*1.52) node[anchor=south] {$B$};
\draw(\x*1.68,0) node[anchor=west] {$C$};
\smalldot{0,0};
\smalldot{\x*0.95,\x*1.52};
\smalldot{\x*1.68,0};
\draw(0,0) -- ++ (96.88:1);
\draw(0,0) -- ++ (96.88- 2*72:1);
\draw(\x*0.95,\x*1.52) -- ++ (252.34+72:1);
\draw(\x*0.95,\x*1.52) -- ++ (252.34-72:1);
\draw(\x*1.68,0) -- ++ (210.61 - 2*72:1);
\draw(\x*1.68,0) -- ++ (210.61:1);
\draw(25:0.65) node {$\arc_A$};
\draw(75:0.7) node {$\theta_{ABC}$};
\draw(-20:0.6) node {$\theta_{ACB}$};
\draw($(\x*0.95,\x*1.52)+(-45:0.7)$) node {$\theta_{BCA}$};
\draw($(\x*0.95,\x*1.52)+(-95:0.65)$) node {$\arc_B$};
\draw($(\x*0.95,\x*1.52)+(210:0.55)$) node {$\theta_{BAC}$};
\draw($(\x*1.68,0)+(90:0.65)$) node {$\theta_{CBA}$};
\draw($(\x*1.68,0)+(140:0.60)$) node {$\arc_C$};
\draw($(\x*1.68,0)+(195:0.6)$) node {$\theta_{CAB}$};
\end{scope}
}

\clearpage
\def\currentfigurenumber{A.5}
% Exact TikZ call from pentagon_v005.tex; no geometric modification.
\tikzfig{cord-delta}{Coordinates for $\Delta$-types}
{
[scale=1.0]
\threepentnoD
{0.00}{0.00}{-5.16}%
{0.99}{1.70}{235.38}%
{1.98}{0.00}{183.43}; %
\draw (0,0) node {$C$};
\draw (0.99,1.70) node {$B$};
\draw (1.98,0.0) node {$A$};
\draw (1.0,0.6) node {$\Delta$};
\draw (-5.16:1) -- ++ (126 - 5.16 :2.5) node[anchor=south] {$\alpha$};
\draw (-5.16:1) -- ++ (- 180 + 126 - 5.16 :1.5) node[anchor=west] {$\beta$};
\draw ++ (1.98,0) ++ (183.43 - 72:1) 
-- ++ (54 + 3.43:1) node[anchor=south] {$\gamma$};
\draw (72-5.16:1) -- ++ (126 + 90 - 5.16:0.5);
\draw (0.99,1.70) ++ (235.38:1) node[anchor = south west] {$x_\alpha$} 
-- ++ (126+90 - 5.16:0.5);
}

\clearpage
\def\currentfigurenumber{A.6}
% Exact TikZ call from pentagon_v005.tex; no geometric modification.
\tikzfig{cord-pinwheel}{Coordinates for pinwheel type}{
\begin{scope}[xshift=4cm,scale=1.2]
\threepentnoD{0.00}{0.00}{46.69}%
{0.82}{1.53}{218.09}%
{1.73}{0.00}{163.28};%
\draw(0,0) node {$A$};
\draw(0.82,1.53) node {$C$};
\draw(1.73,0) node {$B$};
\draw(1.8,1.2) node {$\alpha$};
\draw(-0.2,1.0) node {$\beta$};
\draw(0.85,-0.7) node {$\gamma$};
\draw(0.9,0.5) node {$P$};
\draw(0.5,0.5) node {$x_\gamma$};
\smalldot{0.772,0.288}; %
\smalldot{0.686,0.728}; %
\end{scope}
}

\clearpage
\def\currentfigurenumber{A.7}
% Exact TikZ call from pentagon_v005.tex; no geometric modification.
\tikzfig{cord-L}{Coordinates for $LJ$-junction type}{
\begin{scope}[xshift=8cm,scale=1.2]
\threepentnoD{0.00}{0.00}{80.86}
{0.97}{1.58}{232.21}
{1.82}{0.00}{223.24};
\draw (0,0) node {$C$};
\draw (0.97,1.58) node {$B$};
\draw (1.82,0) node {$A$};
\draw (2.0,1.2) node {$\gamma$};
\draw (-0.1,1.2) node {$\alpha$};
\draw (0.85,-0.8) node {$\beta$};
\draw (0.85,0.75) node {$L$};
\draw (0.55,0.4) node {$x_\alpha$};
\smalldot{$(0,0)+(80.86-72:1)$};
\smalldot{$(0.97,1.58)+(232.24:1)$};
\end{scope}
}

\clearpage
\def\currentfigurenumber{A.8}
% Exact TikZ call from pentagon_v005.tex; no geometric modification.
\tikzfig{cord-T}{Coordinates for $TJ$-junction types}{
\begin{scope}[xshift=12cm,scale=1.2]
\threepentnoD{0.00}{0.00}{114.48} %
{0.90}{1.59}{237.18} %
{1.66}{0.00}{219.24}; %
\draw (0,0) node {$A$};
\draw (0.9,1.59) node {$B$};
\draw (1.66,0) node {$C$};
\draw (0.9,0.9) node {$P$};
\draw (0,1.1) node {$\gamma$};
\draw (0.75,-0.9) node {$\beta$};
\draw (1.9,1.2) node {$\alpha$};
\draw (0.45,0.5) node {$x_\gamma$};
\smalldot{0.737,0.675};
\smalldot{0.358,0.75};
\end{scope}
}

\clearpage
\def\currentfigurenumber{A.9}
% Exact TikZ call from pentagon_v005.tex; no geometric modification.
\tikzfig{cord-pint}{Coordinates for pin-$T$ junction types. Although
  it is difficult to tell from the figure, $A$ points into $B$, $B$
  into $C$, and $C$ into $A$.  The parameter $\beta'\approx 0$
  measures the incidence angle between the nearly horizontal edges of
  $A$ and $C$.  The region bounded by the three pentagons is a
  $T$-shaped pentagon, with stem between the nearly parallel edges of
  $A$ and $C$ and two arms along $B$.  The arm between $B$ and $C$ is
  almost imperceptible.  The distance between the neighboring vertices
  of $B$ and $C$ is $x_\alpha$.  The figure on the right distorts the
  pentagons to make the incidence relations more apparent.}
{
\begin{scope}[xshift=4.5cm,yshift=1.5cm]
\threepentnoD{0.00}{0.00}{90} %
{1.40}{-1.377}{0} %
{-0.35}{-1.628}{-18.89}; %
\draw (0,0) node {$C$};
\draw (1.4,-1.37) node {$B$};
\draw (-0.35,-1.62) node {$A$};
\draw (0.36,-2.4) node {$\gamma$};
\draw (-1.1,-0.5) node {$\beta$};
\draw (1.0,-0.3) node {$\alpha$};
\draw (0.4,-0.49) node {$x_\alpha$};
\draw (1.4 - 0.809,-1.377 - 0.5878) -- +(0,-0.75);
\end{scope}
\begin{scope}[xshift=9.5cm,yshift=1cm]
\draw[red] (0.2,-1.2) -- ++ (90:1.4) -- ++ (10:1); %
\draw[red] (0.2,-1.2) -- ++ (-10:1);
\draw[red] (-1,0) -- ++(1,0) -- ++ (45:1); %
\draw[red] (-1,0) -- ++(100:1);
\draw[red] (-1,0) -- + (-10:0.9) -- (0.2,-0.8) -- ++ (-120:1); %
\draw[red] (-1,0) -- ++(180-10:0.3) -- ++(-110:1);
\draw (0.9,-0.4) node {$B$};
\draw (-0.5,0.7) node {$C$};
\draw (-0.7,-0.7) node {$A$};
\end{scope}
}

\clearpage
\bibliographystyle{plainnat}
\bibliography{../paper/pentagon}
\end{document}
